% =============================================================================
%  Cup Manipulator + SEA Cable: Full System Description & C3M Applicability
%  Complete state-space derivation for contraction-based control design
% =============================================================================

\documentclass[11pt, a4paper]{article}

\usepackage{amsmath, amssymb, bm}
\usepackage{geometry}
\usepackage[pdfusetitle, bookmarksnumbered]{hyperref}
\usepackage{booktabs}
\usepackage{xcolor}
\usepackage{cleveref}
\usepackage[backgroundcolor=blue!5, linecolor=blue!40, linewidth=1pt,
            innertopmargin=6pt, innerbottommargin=6pt,
            innerleftmargin=7pt, innerrightmargin=7pt]{mdframed}

\geometry{margin=2.2cm}

% ── Standard project macros ──────────────────────────────────────────────────
\newcommand{\bq}{\bm{q}}
\newcommand{\bp}{\bm{p}}
\newcommand{\bu}{\bm{u}}
\newcommand{\bx}{\bm{x}}
\newcommand{\btau}{\bm{\tau}}
\newcommand{\bF}{\bm{F}}
\newcommand{\bJ}{\bm{J}}
\newcommand{\bM}{\bm{M}}
\newcommand{\bC}{\bm{C}}
\newcommand{\bG}{\bm{G}}
\newcommand{\bB}{\bm{B}}
\newcommand{\bW}{\bm{W}}
\newcommand{\R}{\mathbb{R}}

\title{\textbf{Cup Manipulator + SEA Cable System}\\[4pt]
       \textbf{Full State-Space Derivation \& C3M Applicability}\\[6pt]
       \large \texttt{script\_cup\_manipulator\_pendulam\_tendon\_with\_spring\_sea\_pydrake.py}}
\date{April 2026}
\author{Isaac-sim robotics notes}

\begin{document}
\maketitle
\tableofcontents
\newpage

% ═════════════════════════════════════════════════════════════════════════════
\section{System Overview}
% ═════════════════════════════════════════════════════════════════════════════

The cup manipulator is a planar 2-DOF serial arm with cable-driven actuation
on joint~2 via a Series Elastic Actuator (SEA).  The physical topology is:

\begin{mdframed}
\textbf{Joint 1 (shoulder, \texttt{link1\_base}):}
Motor shaft $\to$ gearbox $\to$ rigid link --- \emph{no compliance}.
Torque $\tau_1$ is applied directly.

\medskip
\textbf{Joint 2 (elbow, \texttt{link2\_link1}):}
Motor drum $\to$ cable $\to$ big pulley $\to$ \textbf{spring} $\to$ link~2.
The spring (belt elasticity or explicit compliance element) is ``in series''
between the motor and the joint.
\end{mdframed}

This document derives the complete 6D state-space model in control-affine form
$\dot{\bx} = \bm{f}(\bx) + \bB(\bx)\,\bu$
and discusses how C3M (Control Contraction Metrics) can be applied.


% ═════════════════════════════════════════════════════════════════════════════
\section{Physical Parameters}
\label{sec:params}
% ═════════════════════════════════════════════════════════════════════════════

\begin{table}[h]
\centering
\caption{Manipulator link parameters (from Onshape URDF export).}
\label{tab:links}
\begin{tabular}{@{}lcrl@{}}
\toprule
Symbol & Description & Value & Unit \\
\midrule
$L_1$   & Link 1 length (shoulder $\to$ elbow) & 334.80 & mm \\
$L_2$   & Link 2 length (elbow $\to$ EE)       & 190.00 & mm \\
$m_1$   & Link 1 mass                          & 5.313  & kg \\
$m_2$   & Link 2 mass                          & 0.520  & kg \\
$r_p$   & Big pulley pitch radius (HTD 5M 60T) & 47.75  & mm \\
\bottomrule
\end{tabular}
\end{table}

\begin{table}[h]
\centering
\caption{Motor parameters (CubeMars AK60-6 KV80).}
\label{tab:motor}
\begin{tabular}{@{}lcrl@{}}
\toprule
Symbol & Description & Value & Unit \\
\midrule
$N$         & Gear ratio                  & 6.0       & -- \\
$J_m$       & Rotor inertia (motor-side)  & $3.32 \times 10^{-5}$ & kg\,m$^2$ \\
$b_m$       & Motor viscous damping       & $0.12/N^2 = 3.33 \times 10^{-3}$ & Nm\,s/rad \\
$\tau_\text{peak}$ & Peak joint torque    & 9.0       & Nm \\
$\tau_\text{cont}$ & Continuous joint torque & 4.5    & Nm \\
\bottomrule
\end{tabular}
\end{table}

\begin{table}[h]
\centering
\caption{SEA cable parameters (typical defaults).}
\label{tab:sea}
\begin{tabular}{@{}lcrl@{}}
\toprule
Symbol & Description & Default & Unit \\
\midrule
$k_s$   & Cable spring stiffness & 30--5000 & N/m \\
$b_c$   & Cable dashpot damping  & 2.0      & N\,s/m \\
\bottomrule
\end{tabular}
\end{table}


% ═════════════════════════════════════════════════════════════════════════════
\section{Subsystem 1: Rigid 2R Manipulator Dynamics}
\label{sec:manip}
% ═════════════════════════════════════════════════════════════════════════════

\subsection{Generalized Coordinates}

Let $\bq = [q_1, q_2]^\top$ denote the joint angles, where $q_1$ is the
shoulder (link1\_base) and $q_2$ is the elbow (link2\_link1).

The standard manipulator equation of motion is:
\begin{equation}
\label{eq:manip-eom}
\bM(\bq)\,\ddot{\bq} + \bC(\bq, \dot{\bq})\,\dot{\bq} + \bG(\bq) = \btau_\text{applied}
\end{equation}
where $\bM \in \R^{2 \times 2}$ is the mass matrix,
$\bC \in \R^{2 \times 2}$ is the Coriolis matrix,
$\bG \in \R^2$ is the gravity vector, and
$\btau_\text{applied} \in \R^2$ is the vector of applied joint torques.

\subsection{Mass Matrix}

For a planar 2R arm with link masses $m_1, m_2$, lengths $L_1, L_2$,
and centres of mass at distances $l_{c1}, l_{c2}$ along each link:
\begin{equation}
\label{eq:mass-matrix}
\bM(\bq) =
\begin{bmatrix}
I_1 + I_2 + m_2 L_1^2 + 2 m_2 L_1 l_{c2} \cos q_2
  & I_2 + m_2 L_1 l_{c2} \cos q_2 \\
I_2 + m_2 L_1 l_{c2} \cos q_2
  & I_2
\end{bmatrix}
\end{equation}
where $I_i = m_i l_{ci}^2$ are the link inertias about their pivot.


\subsection{Applied Torques}

\begin{mdframed}
\textbf{Crucially}, $\btau_\text{applied}$ is \emph{not} the controller output.
\begin{itemize}
\item \textbf{Joint 1} (rigid): $\tau_{1,\text{applied}} = \tau_1$ \quad (direct drive, pass-through)
\item \textbf{Joint 2} (SEA):   $\tau_{2,\text{applied}} = r_p \cdot F_\text{cable}$
      \quad (spring force, \emph{not} $\tau_{2,\text{des}}$)
\end{itemize}
The controller commands $[\tau_1,\, \tau_{2,\text{des}}]$, but
$\tau_{2,\text{des}}$ only reaches $q_2$ through the motor and spring dynamics.
\end{mdframed}


% ═════════════════════════════════════════════════════════════════════════════
\section{Subsystem 2: SEA Motor--Spring--Cable Dynamics}
\label{sec:sea}
% ═════════════════════════════════════════════════════════════════════════════

\subsection{Spring Extension}

The cable wraps on the motor drum (motor-side) and the big pulley
(joint-side).  In torque mode, the motor-side angle is $\theta_m$
(at the rotor, before the gearbox).  At the gearbox output, the angle
is $\theta_m / N$.  The linear cable displacement from the motor is:
\[
l_\text{motor} = r_p \cdot \frac{\theta_m}{N}
\]
and from the joint:
\[
l_\text{joint} = r_p \cdot q_2
\]
The \textbf{linear spring extension} is:
\begin{equation}
\label{eq:delta}
\boxed{\delta = r_p \left(\frac{\theta_m}{N} - q_2\right)}
\end{equation}

\subsection{Cable Force (Spring--Damper)}

The spring--damper produces a force:
\begin{equation}
\label{eq:F-raw}
F_\text{raw} = k_s\,\delta + b_c\,\dot{\delta}
\end{equation}
where $\dot{\delta} = r_p \left(\dot{\theta}_m / N - \dot{q}_2\right)$.

\paragraph{Unilateral cable constraint.}
Cables can only \emph{pull} (tension $\ge 0$), never push.
The antagonistic cable pair gives:
\begin{equation}
\label{eq:cable-tensions}
\delta > 0: \quad T_\text{green} = \max(F_\text{raw}, 0),\; T_\text{red} = 0
\qquad
\delta < 0: \quad T_\text{green} = 0,\; T_\text{red} = \max(-F_\text{raw}, 0)
\end{equation}
The net cable force is:
\begin{equation}
\label{eq:F-cable}
F_\text{cable} = T_\text{green} - T_\text{red}
\end{equation}
and the torque applied to joint~2 is $\tau_{2,\text{applied}} = r_p \cdot F_\text{cable}$.

\subsection{Motor Rotor Dynamics (Torque Mode)}

The motor in MIT torque mode obeys 2nd-order rotor dynamics:
\begin{equation}
\label{eq:motor}
\boxed{J_m\,\ddot{\theta}_m = \tau_m - b_m\,\dot{\theta}_m - \frac{r_p\,F_\text{cable}}{N}}
\end{equation}
where the motor-side torque command is:
\[
\tau_m = \frac{\tau_{2,\text{des}}}{N}
\]

The effective torsional spring stiffness seen by the rotor is
$r_p^2 k_s / N^2$, giving a motor-side resonance frequency:
\begin{equation}
\omega_n = \sqrt{\frac{r_p^2\,k_s}{N^2\,J_m}}
\end{equation}
For AK60-6 with $k_s = 300\;\text{N/m}$: $\omega_n \approx 24\;\text{rad/s}$ (4\,Hz).


% ═════════════════════════════════════════════════════════════════════════════
\section{Combined 6D State-Space Model}
\label{sec:statespace}
% ═════════════════════════════════════════════════════════════════════════════

\subsection{State and Control Vectors}

\begin{mdframed}
\textbf{State vector} ($n = 6$):
\begin{equation}
\label{eq:state}
\bx = \begin{bmatrix} q_1 \\ q_2 \\ \dot{q}_1 \\ \dot{q}_2 \\ \theta_m \\ \dot{\theta}_m \end{bmatrix}
\in \R^6
\end{equation}

\textbf{Control input} ($m = 2$):
\begin{equation}
\label{eq:control}
\bu = \begin{bmatrix} \tau_1 \\ \tau_{2,\text{des}} \end{bmatrix}
\in \R^2
\end{equation}
\end{mdframed}

\subsection{Control-Affine Form}

The system can be written as:
\begin{equation}
\label{eq:xdot}
\boxed{\dot{\bx} = \bm{f}(\bx) + \bB(\bx)\,\bu}
\end{equation}

\subsubsection{Drift term \texorpdfstring{$\bm{f}(\bx)$}{f(x)}}

From the manipulator EOM (\ref{eq:manip-eom}) with
$\btau_\text{applied} = [\tau_1, \; r_p F_\text{cable}]^\top$,
and eliminating the control terms:

\begin{equation}
\label{eq:f}
\bm{f}(\bx) =
\begin{bmatrix}
\dot{q}_1   \\[2pt]
\dot{q}_2   \\[2pt]
\left[\bM^{-1}(\bq)\left(-\bC(\bq,\dot{\bq})\dot{\bq} - \bG(\bq)
  + \begin{pmatrix} 0 \\ r_p\,F_\text{cable} \end{pmatrix}\right)\right]_1   \\[4pt]
\left[\bM^{-1}(\bq)\left(-\bC(\bq,\dot{\bq})\dot{\bq} - \bG(\bq)
  + \begin{pmatrix} 0 \\ r_p\,F_\text{cable} \end{pmatrix}\right)\right]_2   \\[4pt]
\dot{\theta}_m   \\[2pt]
\dfrac{1}{J_m}\left(-b_m\,\dot{\theta}_m - \dfrac{r_p\,F_\text{cable}}{N}\right)
\end{bmatrix}
\end{equation}

Note that $F_\text{cable} = F_\text{cable}(\bx)$ depends on the state
through \eqref{eq:delta}--\eqref{eq:F-cable}.
The spring force \emph{couples} the motor state $(\theta_m, \dot{\theta}_m)$
to the joint state $(q_2, \dot{q}_2)$ --- this coupling lives entirely in
$\bm{f}$, \emph{not} in $\bB$.

\subsubsection{Input matrix \texorpdfstring{$\bB(\bx)$}{B(x)}}

Writing $\bM^{-1}(\bq) = \begin{bmatrix} M^{-1}_{11} & M^{-1}_{12} \\
M^{-1}_{21} & M^{-1}_{22} \end{bmatrix}$:

\begin{equation}
\label{eq:B}
\bB(\bx) =
\begin{bmatrix}
0           & 0                \\
0           & 0                \\
M^{-1}_{11}(\bq) & 0          \\
M^{-1}_{21}(\bq) & 0          \\
0           & 0                \\
0           & \dfrac{1}{J_m\,N}
\end{bmatrix}
\in \R^{6 \times 2}
\end{equation}

\begin{mdframed}
\textbf{Key observation}: $\tau_{2,\text{des}}$ appears \emph{only} in
row 6 ($\ddot{\theta}_m$), \emph{not} in row 4 ($\ddot{q}_2$).
The commanded torque reaches joint~2 \emph{only through the spring dynamics}
in $\bm{f}(\bx)$.  This is the fundamental source of torque lag when $k_s$ is
small.
\end{mdframed}

\subsection{Structure of \texorpdfstring{$\bB$}{B}: Row-by-Row}

\begin{table}[h]
\centering
\caption{Input matrix $\bB(\bx)$ structure.}
\label{tab:B-structure}
\begin{tabular}{@{}clcc@{}}
\toprule
Row & State derivative & $\tau_1$ column & $\tau_{2,\text{des}}$ column \\
\midrule
0 & $\dot{q}_1$            & 0    & 0 \\
1 & $\dot{q}_2$            & 0    & 0 \\
2 & $\ddot{q}_1$           & $M^{-1}_{11}(\bq)$ & 0 \\
3 & $\ddot{q}_2$           & $M^{-1}_{21}(\bq)$ & \textbf{0} \\
4 & $\dot{\theta}_m$       & 0    & 0 \\
5 & $\ddot{\theta}_m$      & 0    & $1/(J_m N)$ \\
\bottomrule
\end{tabular}
\end{table}

The input matrix has $\text{rank}(\bB) = 2$ everywhere (since $M^{-1}_{11} > 0$
and $J_m > 0$), but $\tau_{2,\text{des}}$ acts only on the motor rotor.


% ═════════════════════════════════════════════════════════════════════════════
\section{Underactuation Analysis}
\label{sec:underact}
% ═════════════════════════════════════════════════════════════════════════════

\begin{table}[h]
\centering
\caption{Actuation comparison: rigid vs.\ SEA manipulator.}
\label{tab:actuation}
\begin{tabular}{@{}lccc@{}}
\toprule
Property & Rigid 2R & SEA Cable (Torque mode) \\
\midrule
States ($n$)   & 4 & \textbf{6} \\
Controls ($m$) & 2 & 2 \\
Unactuated directions ($n - m$) & 2 & \textbf{4} \\
Fully actuated? & Yes & \textbf{No} \\
\bottomrule
\end{tabular}
\end{table}

\noindent
The rigid manipulator is fully actuated: both $\tau_1$ and $\tau_2$ act
directly on $\ddot{q}_1$ and $\ddot{q}_2$ through the invertible mass
matrix, and feedback linearisation (computed torque) is exact.

With the SEA, $\tau_{2,\text{des}}$ only enters $\ddot{\theta}_m$.
Joint~2 is driven \emph{indirectly} through the spring coupling in $\bm{f}$.
The system is \textbf{underactuated}: we cannot independently command $\ddot{q}_2$.

\subsection{Null Space of \texorpdfstring{$\bB^\top$}{B-transpose}}

The null space $\bB^\perp$ satisfies $\bB^\top \bB^\perp = \bm{0}$.
Since $\bB$ has rank~2 in $\R^6$, the null space has dimension
$n - m = 4$.  A basis is:

\begin{equation}
\label{eq:Bbot}
\bB^\perp(\bx) =
\begin{bmatrix}
1 & 0 & 0 & 0 \\
0 & 1 & 0 & 0 \\
0 & 0 & 0 & 0 \\
0 & -M^{-1}_{21}/M^{-1}_{11} & 1 & 0 \\
0 & 0 & 0 & 1 \\
0 & 0 & 0 & 0
\end{bmatrix}
\in \R^{6 \times 4}
\end{equation}

The columns correspond to:
\begin{enumerate}
\item $\bm{e}_1$: the $q_1$ kinematic identity $(\dot{q}_1 = \dot{q}_1)$
\item Modified $\bm{e}_2$: the $q_2$ direction, with row~4 adjusted so
      the $\tau_1$ coupling through $M^{-1}_{21}$ is cancelled
\item $\bm{e}_4$: the $\ddot{q}_2$ direction decoupled from $\tau_1$
\item $\bm{e}_5$: the $\theta_m$ kinematic identity
\end{enumerate}


% ═════════════════════════════════════════════════════════════════════════════
\section{Current Control: Computed Torque + SEA}
\label{sec:ct}
% ═════════════════════════════════════════════════════════════════════════════

The existing controller in \texttt{controller/controller.py} uses a
\textbf{two-block architecture}:

\[
\text{Trajectory} \xrightarrow{\bp_\text{ref}(t)}
\text{CT Controller} \xrightarrow{[\tau_1,\,\tau_{2,\text{des}}]}
\text{SEA Actuator} \xrightarrow{[\tau_1,\,r_p F_\text{cable}]}
\text{MultibodyPlant}
\]

The CT controller computes:
\begin{align}
\bq_\text{des} &= \text{IK}(\bp_\text{ref}) \label{eq:ik}  \\
\bm{a}_\text{des} &= \ddot{\bq}_\text{ref}
  + K_p\,(\bq_\text{des} - \bq)
  + K_d\,(\dot{\bq}_\text{ref} - \dot{\bq}) \label{eq:pd}  \\
\btau_\text{des} &= \bM(\bq)\,\bm{a}_\text{des}
  + \bC(\bq, \dot{\bq})\,\dot{\bq}
  + \bG(\bq) \label{eq:invdyn}
\end{align}

\begin{mdframed}
\textbf{Why CT fails with soft springs:}
The CT law \eqref{eq:invdyn} assumes $\btau_\text{applied} = \btau_\text{des}$
(instant torque authority).  With the SEA, $\tau_{2,\text{applied}} = r_p\,F_\text{cable}
\neq \tau_{2,\text{des}}$, and the mismatch grows as $k_s$ decreases.
The CT controller is \emph{model-blind} to the spring dynamics --- it treats
the SEA as if it were rigid.
\end{mdframed}


% ═════════════════════════════════════════════════════════════════════════════
\section{C3M: Control Contraction Metrics}
\label{sec:c3m}
% ═════════════════════════════════════════════════════════════════════════════

\subsection{Background}

C3M (Sun et al., 2020) jointly learns two neural networks:
\begin{enumerate}
\item $\bW(\bx)$: a state-dependent \textbf{dual contraction metric}
      (positive definite matrix)
\item $\bu(\bx, \bx_e, \bu_\text{ref})$: a \textbf{tracking controller}
      that guarantees exponential convergence at rate $\lambda$
\end{enumerate}
where $\bx_e = \bx - \bx_\text{ref}$ is the tracking error and
$\bu_\text{ref}$ is the feedforward control.

\subsection{Contraction Condition}

With $\bm{M} = \bW^{-1}$ (the contraction metric), the training loss enforces:
\begin{equation}
\label{eq:contraction}
\dot{\bm{M}} + (\bm{A} + \bB\bm{K})^\top \bm{M}
  + \bm{M}\,(\bm{A} + \bB\bm{K})
  + 2\lambda\,\bm{M} \preceq 0
\end{equation}
where $\bm{A} = \partial\bm{f}/\partial\bx + \sum_i u_i\,\partial\bB_i/\partial\bx$
and $\bm{K} = \partial\bu/\partial\bx$.  This ensures the distance between
any two trajectories (measured by $\bm{M}$) contracts exponentially at rate~$\lambda$.

\subsection{Dual Conditions (C1 and C2)}

C3M decomposes the contraction condition into two parts projected through
$\bB^\perp$:

\paragraph{C1 --- Unactuated contraction.}
\begin{equation}
\label{eq:C1}
(\bB^\perp)^\top
\left[
-\dot{\bW}\bm{f}
+ \frac{\partial\bm{f}}{\partial\bx}\bW
+ \bW\left(\frac{\partial\bm{f}}{\partial\bx}\right)^\top
+ 2\lambda\,\bW
\right]
\bB^\perp \preceq 0
\end{equation}
This must hold in the \textbf{4-dimensional null space} of $\bB^\top$.
For the SEA system, this subspace includes the $q_2$ and $\theta_m$ directions
--- so C3M must find a metric that certifies the \emph{spring dynamics naturally
contract} even though we cannot directly command $\ddot{q}_2$.

\paragraph{C2 --- Actuated orthogonality.}
\begin{equation}
\label{eq:C2}
(\bB^\perp)^\top
\left[
\dot{\bW}\bB_j
- \frac{\partial\bB_j}{\partial\bx}\bW
- \bW\left(\frac{\partial\bB_j}{\partial\bx}\right)^\top
\right]
\bB^\perp = 0
\quad \forall\; j = 1, \ldots, m
\end{equation}
This ensures the control directions are metrically orthogonal to the
unactuated subspace.

\subsection{Why C3M is Well-Suited for This System}

\begin{table}[h]
\centering
\caption{Comparison of control approaches for the cup manipulator + SEA.}
\label{tab:comparison}
\begin{tabular}{@{}lcccc@{}}
\toprule
Property & CT & CT + RL residual & LQR & \textbf{C3M} \\
\midrule
Accounts for spring dynamics & No & Learned & Linear only & \textbf{Yes} \\
Guaranteed convergence rate  & No & No & Local only & \textbf{Yes ($\lambda$)} \\
Handles underactuation       & No & Partially & Partially & \textbf{Yes} \\
Robust to model uncertainty  & No & Partially & No & \textbf{Certified} \\
Need training                & No & Yes & No & Yes \\
\bottomrule
\end{tabular}
\end{table}

\noindent
Key advantages of C3M for this specific system:

\begin{enumerate}
\item \textbf{Spring dynamics awareness.}
      The C1 condition \eqref{eq:C1} explicitly verifies that the spring
      coupling from $\theta_m$ to $q_2$ is contracting.  If $k_s$ is too
      small, C1 becomes infeasible for a given $\lambda$ --- this gives a
      \emph{principled lower bound on spring stiffness}.

\item \textbf{Anticipatory control.}
      Unlike CT (which ignores the spring lag), C3M learns a controller
      that \emph{anticipates} the motor-spring dynamics.  The controller
      $\bu(\bx, \bx_e, \bu_\text{ref})$ sees the full state including
      $(\theta_m, \dot{\theta}_m)$ and can pre-compensate for the spring lag.

\item \textbf{Convergence rate certificate.}
      The parameter $\lambda > 0$ is a guaranteed convergence rate.
      For tracking tasks, this translates to a bound on the steady-state
      tracking error: fast trajectories with low $\lambda$ will have larger
      error, providing a principled speed--accuracy tradeoff.

\item \textbf{Unilateral cable handling.}
      The $\max(F, 0)$ nonlinearity can be smoothed via softplus
      $\tilde{F} = \frac{1}{\beta}\ln(1 + e^{\beta F})$ during training
      to maintain differentiability for PyTorch autograd, while the
      deployment uses exact $\max$.
\end{enumerate}


% ═════════════════════════════════════════════════════════════════════════════
\section{C3M Integration Architecture}
\label{sec:integration}
% ═════════════════════════════════════════════════════════════════════════════

\subsection{Training Phase}

Three files must be created in the C3M framework
(\texttt{contraction-theory/C3M/}):

\begin{enumerate}
\item \texttt{systems/system\_CUPMANIP\_SEA.py} --- 6D $\bm{f}(\bx)$,
      $\bB(\bx)$, $\bB^\perp(\bx)$ in PyTorch
\item \texttt{configs/config\_CUPMANIP\_SEA.py} --- state/control bounds,
      sampling ranges
\item \texttt{models/model\_CUPMANIP\_SEA.py} --- neural network
      architecture for $\bW$ and $\bu$
\end{enumerate}

The $\bm{f}(\bx)$ implementation requires symbolic 2R manipulator EOM
in PyTorch (not Drake), since autograd must compute $\partial\bm{f}/\partial\bx$.

Training command:
\begin{verbatim}
cd contraction-theory/C3M
python main.py --task CUPMANIP_SEA --lambda 1.5 \
    --log checkpoints/cupmanip_sea/lambda_1.5 --epochs 15
\end{verbatim}

\subsection{Deployment in Drake}

The trained C3M controller replaces the CT controller as a Drake
\texttt{LeafSystem}:

\[
\text{Trajectory}
\xrightarrow{\bx_\text{ref}(t)}
\texttt{C3MController}
\xrightarrow{[\tau_1,\,\tau_{2,\text{des}}]}
\underbrace{\texttt{SEACableActuator}}_{\text{physical spring model}}
\xrightarrow{[\tau_1,\,r_p F]}
\texttt{Plant}
\]

The SEA actuator block \emph{remains unchanged} --- it models physical reality.
The C3M controller produces torque commands that account for the spring dynamics,
rather than blindly inverting the manipulator dynamics and hoping the spring
keeps up.

The reference state includes the motor:
\begin{equation}
\bx_\text{ref}(t) = \begin{bmatrix}
q_{1,\text{ref}}(t) \\
q_{2,\text{ref}}(t) \\
\dot{q}_{1,\text{ref}}(t) \\
\dot{q}_{2,\text{ref}}(t) \\
N \cdot q_{2,\text{ref}}(t) \\
N \cdot \dot{q}_{2,\text{ref}}(t)
\end{bmatrix}
\end{equation}
where $\theta_{m,\text{ref}} = N \cdot q_{2,\text{ref}}$ corresponds to
the spring-at-rest equilibrium ($\delta = 0$).

\subsection{Neural Network Architecture}

The controller network in C3M takes an \texttt{effective\_dim} slice of the
state as input.  For the cart-pendulum (4D state, 1D control), dimensions
1--3 ($\dot{p}, \theta, \dot{\theta}$) are used (cart position $p$ is ignored
since the dynamics are translation-invariant).

For the cup manipulator SEA (6D state, 2D control), \textbf{all 6 states
matter} because the dynamics depend on all angles and the motor state:
\[
\texttt{effective\_dim\_start} = 0, \quad
\texttt{effective\_dim\_end} = 6
\]

The network dimensions scale accordingly:
\begin{itemize}
\item $\bW$: \texttt{Linear}$(6, 128) \to \tanh \to$ \texttt{Linear}$(128, 36)$
      \quad ($6 \times 6 = 36$ entries)
\item $\bu_{w1}$: \texttt{Linear}$(12, 128) \to \tanh \to$ \texttt{Linear}$(128, 6c)$
      \quad (input: $[\bx_\text{eff},\, \bx_e{}_\text{eff}]$, $c = 3n = 18$)
\item $\bu_{w2}$: \texttt{Linear}$(12, 128) \to \tanh \to$ \texttt{Linear}$(128, 2c)$
      \quad (output: $2 \times c$ matrix)
\end{itemize}


% ═════════════════════════════════════════════════════════════════════════════
\section{Recommended Development Path}
\label{sec:roadmap}
% ═════════════════════════════════════════════════════════════════════════════

\begin{enumerate}
\item \textbf{Phase 1: Rigid 4D system.}
      Start with the rigid manipulator ($n = 4, m = 2$, fully actuated).
      This should train easily and validate the PyTorch dynamics
      against Drake.  $\bB^\perp$ is only $4 \times 2$
      (two kinematic identities).

\item \textbf{Phase 2: 6D SEA system.}
      Add the motor states.  The 4D null space and spring nonlinearity
      make training harder.  Use softplus approximation for $\max(F, 0)$.
      Start with stiff springs ($k_s = 500$) and gradually reduce.

\item \textbf{Phase 3: Feasibility boundary.}
      Sweep $\lambda$ vs.\ $k_s$ to find the region where C1 is
      satisfiable.  This gives a principled minimum-stiffness curve
      for a given convergence rate --- useful for hardware design.

\item \textbf{Phase 4: Drake deployment.}
      Package the trained C3M controller as a Drake \texttt{LeafSystem},
      compare against CT + SEA and CT + RL residual on the same
      rectangular tracking trajectory.
\end{enumerate}

\end{document}
